\documentclass[a4paper,12pt]{report}
\usepackage[spanish]{babel}
\usepackage[utf8]{inputenc}
\usepackage{graphicx}
\usepackage{amsmath}
\usepackage{hyperref}
\usepackage{listings}
\usepackage{float}
\usepackage{booktabs}
\usepackage{geometry}
\usepackage{tabularx}

% Configuración de listings para código
\lstdefinestyle{pythonstyle}{
    language=Python,
    basicstyle=\ttfamily\footnotesize,
    breaklines=true,
    numbers=left,
    numberstyle=\tiny,
    showstringspaces=false,
    frame=single
}

\geometry{margin=2.5cm}
\hypersetup{
  colorlinks=true,
  linkcolor=blue,
  urlcolor=blue,
  citecolor=blue
}

\title{\textbf{Mapeando el Caos: Indexación y Organización de Datos No Estructurados para Datos Multimedia}\\[1ex]
\large CS2702 – Base de Datos 2\\
\large Universidad de Ingeniería y Tecnología (UTEC)}
\author{
  Ian Kevin Condori Flores \\
  Angel Tribeño Hurtado \\
  Nicolás Stigler Yañez \\
  Rafael Humberto Ramos Huamaní \\
}
\date{\today}
\begin{document}

\maketitle
\newpage
\tableofcontents
\newpage

\chapter{Introducción}
\section{Contexto}
En este proyecto se aborda la indexación y búsqueda en colecciones multimodales, incluyendo texto, audio e imágenes. El objetivo es diseñar una base de datos que permita consultas SQL extendidas (por ejemplo, operadores de similitud) y gestionarlas mediante un único backend en Python.
Se utiliza:
\begin{itemize}
  \item Un índice invertido textual (clase \texttt{InvertedIndex} en \texttt{index/inverted\_index.py}) para recuperar documentos basados en similitud coseno.
  \item Motores de búsqueda multimedia (clases \texttt{MultimediaSequentialIndex} y \texttt{MultimediaInvertedIndex} en \texttt{indexes/multimediatree.py}) para imágenes y audio.
  \item Un parser SQL propio (módulo \texttt{parser/parser.py}) que normaliza consultas y despacha al motor correspondiente.
\end{itemize}
La importancia de una base de datos multimodal radica en la capacidad de realizar búsquedas por contenido, más allá de texto estructurado, integrando metadatos y similitud perceptual.

\chapter{Construcción del Índice Invertido Textual}
\section{Preprocesamiento}
Antes de indexar, el texto se normaliza, tokeniza, elimina \emph{stopwords} y se aplica stemming utilizando el módulo \texttt{backend/utils/text\_processing.py}.

El proceso incluye:
\begin{itemize}
\item Conversión a minúsculas
\item Tokenización usando NLTK
\item Eliminación de stopwords
\item Aplicación de stemming con SnowballStemmer
\end{itemize}

\section{Construcción del Índice}
La clase \texttt{InvertedIndex} agrega documentos con TF, calcula IDF al finalizar y almacena pesos TF-IDF.

El proceso de construcción incluye:
\begin{enumerate}
\item Cálculo de frecuencia de términos (TF)
\item Actualización de normas de documentos
\item Cálculo de IDF al finalizar
\item Persistencia del índice en disco
\end{enumerate}

\section{Consulta}
Para recuperar los documentos más similares se implementa \texttt{search\_knn}, que crea un vector de consulta TF-IDF y calcula similitud coseno con todos los documentos indexados.

\chapter{Indexación de Descriptores Locales en Multimedia}
\section{Procesamiento y Búsqueda en Audio}
\subsection{Extracción de Características de Audio}
Para el procesamiento de audio se utilizan coeficientes MFCC (Mel-Frequency Cepstral Coefficients) que capturan las características espectrales del audio.

El proceso incluye:
\begin{itemize}
\item Carga del archivo de audio usando librosa
\item Extracción de 13 coeficientes MFCC por frame
\item Normalización de características
\end{itemize}

\subsection{Generación de Histogramas TF-IDF para Audio}
Se implementa un enfoque Bag-of-Words para audio:
\begin{enumerate}
\item Construcción de codebook usando K-means
\item Cuantización de descriptores MFCC
\item Generación de histogramas de palabras visuales
\item Aplicación de pesos TF-IDF
\end{enumerate}

\subsection{Indexación y Búsqueda por Similitud en Audio}
La clase \texttt{MultimediaSequentialIndex} permite búsqueda secuencial mientras que \texttt{MultimediaInvertedIndex} utiliza un índice invertido para mayor eficiencia.

\section{Procesamiento y Búsqueda en Imágenes}
\subsection{Extracción de Características SIFT}
Para imágenes se utilizan descriptores SIFT (Scale-Invariant Feature Transform) que son robustos a cambios de escala, rotación e iluminación.

\subsection{Indexación y Búsqueda por Similitud en Imágenes}
Similar al audio, se implementa un enfoque Bag-of-Words para imágenes usando descriptores SIFT cuantizados en un codebook.

\chapter{Índice Invertido Textual en Memoria Secundaria}
\section{Almacenamiento en Disco}
\begin{itemize}
  \item \texttt{inverted\_index.pkl}
  \item \texttt{inverted\_index.txt}
\end{itemize}

\section{Similitud de Coseno y Consultas Eficientes}
\begin{enumerate}
  \item Preprocesamiento de la consulta con \texttt{preprocess\_text}
  \item Cálculo del vector de consulta y su norma
  \item Recorrido de postings
  \item División por normas precomputadas
  \item Selección de top-K
\end{enumerate}

\chapter{Frontend}
\section{Mini-manual de Usuario}
\subsection{Instalación y Configuración}
\begin{enumerate}
\item Instalar dependencias: \texttt{pip install -r requirements.txt}
\item Ejecutar el backend: \texttt{uvicorn raf.backend:app --reload --port 8000}
\item Ejecutar el frontend: \texttt{streamlit run raf/streamlit\_app.py}
\item Navegar a \texttt{http://localhost:8501}
\end{enumerate}

\subsection{Funcionalidades Disponibles}
\subsubsection{Modo SQL}
Permite ejecutar consultas SQL estándar y extendidas:
\begin{itemize}
\item CREATE, INSERT, SELECT, DELETE
\item Soporte para tipos multimedia (AUDIO, IMAGE, TEXT)
\item Consultas de similitud usando operador \texttt{<->}
\end{itemize}

\subsubsection{Modo Texto}
Interface para búsqueda textual:
\begin{itemize}
\item Subida de documentos de texto
\item Construcción de índice invertido
\item Búsqueda por similitud coseno
\end{itemize}

\subsubsection{Modo Multimedia}
Interface para búsqueda multimedia:
\begin{itemize}
\item Selección de archivos locales o subida
\item Construcción de codebook multimedia
\item Búsqueda KNN por similitud
\end{itemize}

\chapter{Experimentación}
\section{Configuración del Experimento}
Se realizaron pruebas de rendimiento en diferentes tipos de consultas:
\begin{itemize}
\item Búsqueda textual con índice invertido
\item Búsqueda multimedia secuencial
\item Búsqueda multimedia con índice invertido
\end{itemize}

\section{Presentación de Resultados}
\begin{table}[H]
\centering
\begin{tabular}{lcc}
\toprule
Tipo de Consulta & Tiempo Promedio & Top-K \\
\midrule
Textual (coseno) & 120 ms & 10 \\
Multimedia Secuencial & 450 ms & 5 \\
Multimedia Invertido & 80 ms & 5 \\
\bottomrule
\end{tabular}
\caption{Comparativa de tiempos de búsqueda}
\end{table}

\section{Análisis de Resultados}
Los resultados muestran que:
\begin{itemize}
\item El índice invertido multimedia es significativamente más rápido que la búsqueda secuencial
\item La búsqueda textual tiene rendimiento intermedio
\item Los tiempos escalan linealmente con el tamaño del dataset
\end{itemize}

\chapter{Conclusiones}
\section{Logros Principales}
\begin{itemize}
\item Implementación exitosa de sistema multimodal
\item Integración de búsqueda textual y multimedia
\item Interface unificada con Streamlit
\item Parser SQL extendido con operadores de similitud
\end{itemize}

\section{Trabajo Futuro}
\begin{itemize}
\item Optimización de algoritmos de extracción de características
\item Implementación de índices más sofisticados (LSH, etc.)
\item Soporte para más tipos de datos multimedia
\item Mejoras en la interface de usuario
\end{itemize}

\end{document}
